% !TEX root = ../main.tex
\begin{figure}[H]
  \centering
  \resizebox{\textwidth}{!}{%
    \begin{tikzpicture}[
      pantalla/.style={draw=diagNodoBorde, thick, rounded corners=4pt, fill=white},
      barra/.style={draw=diagGrupoBorde, fill=diagGrupo, rounded corners=2pt,
                    align=center, font=\footnotesize\bfseries, minimum height=0.8cm},
      usuario/.style={draw=diagNodoBorde, fill=diagNodo, rounded corners=4pt,
                      align=left, font=\scriptsize, inner sep=5pt},
      asistente/.style={draw=diagNodoBorde, fill=white, rounded corners=4pt,
                        align=left, font=\scriptsize, inner sep=5pt},
      pie/.style={font=\tiny, text=black!55, align=left, inner sep=1pt},
      nota/.style={draw=diagGrupoBorde, fill=diagGrupo, rounded corners=3pt,
                   font=\scriptsize, text=black!75, align=left, inner sep=4pt},
      entrada/.style={draw=diagNodoBorde, rounded corners=4pt, fill=white,
                      align=left, font=\scriptsize\itshape, inner sep=5pt},
      punto/.style={circle, draw=diagNodoBorde, fill=diagNodo, minimum size=3.2mm,
                    inner sep=0pt},
      rotulo/.style={font=\small\bfseries, anchor=north},
      ]

      \node[pantalla, minimum width=7.6cm, minimum height=11cm] at (3.8,5.5) {};
      \node[barra, minimum width=7cm] at (3.8,10.3) {Asistente de titulaciones · EPSJ};
      \node[usuario, text width=4.4cm, anchor=north east] at (7.4,9.6)
      {Me gustan los videojuegos y la programación, ¿qué me recomiendas?};
      \node[asistente, text width=5.8cm, anchor=north west] (r1) at (0.3,7.9)
      {¡Genial! Si te gustan los videojuegos y la programación, te recomiendo el
        Grado en Ingeniería Informática. Dentro de este grado, puedes
        especializarte en la mención «Sistemas gráficos» y cursar la asignatura
        optativa de «Desarrollo de videojuegos» (6 ECTS)\ldots};
      \node[pie, anchor=north west, xshift=0.1cm, yshift=-0.1cm] at (r1.south west) {09:01 \quad 82,5 s};
      \node[entrada, text width=6.5cm, anchor=north west] at (0.3,1.5)
      {Escribe tu pregunta\ldots};
      \node[rotulo] at (3.8,-0.2) {1 · Respuesta elaborada};

      \begin{scope}[shift={(8.6,0)}]
        \node[pantalla, minimum width=7.6cm, minimum height=11cm] at (3.8,5.5) {};
        \node[barra, minimum width=7cm] at (3.8,10.3) {Asistente de titulaciones · EPSJ};
        \node[usuario, text width=4.4cm, anchor=north east] at (7.4,9.6)
        {¿Y qué asignaturas tiene en primero?};
        \node[asistente, text width=5.8cm, minimum height=1.9cm, anchor=north west] (r2) at (0.3,8.3) {};
        \node[punto] at (1.0,7.85) {};
        \node[punto] at (1.5,7.85) {};
        \node[punto] at (2.0,7.85) {};
        \node[font=\scriptsize, text width=5.4cm, align=left, anchor=north west] at (0.5,7.5)
        {Redactando la respuesta\ldots\ 47 s. El modelo se ejecuta en este mismo
          equipo, así que tarda más que un servicio en la nube.};
        \node[nota, text width=5.6cm, anchor=north west] at (0.4,5.4)
        {La cuenta empieza en «Buscando en la información de la Escuela» y cambia
          a este texto cuando llegan las fuentes, no al cabo de un tiempo fijo.
          Mediana medida: 62,7 s; percentil 90: 125,3 s.};
        \node[entrada, text width=6.5cm, anchor=north west] at (0.3,1.5)
        {Escribe tu pregunta\ldots};
        \node[rotulo] at (3.8,-0.2) {2 · La espera};
      \end{scope}

      \begin{scope}[shift={(17.2,0)}]
        \node[pantalla, minimum width=7.6cm, minimum height=11cm] at (3.8,5.5) {};
        \node[barra, minimum width=7cm] at (3.8,10.3) {Asistente de titulaciones · EPSJ};
        \node[asistente, text width=5.8cm, anchor=north west] (r3) at (0.3,9.6)
        {¡Hola! Te puedo ayudar con las titulaciones de la Escuela Politécnica
          Superior de Jaén: qué grados y dobles grados se estudian allí\ldots};
        \node[pie, anchor=north west, xshift=0.1cm, yshift=-0.1cm] at (r3.south west)
        {09:00 \quad respuesta inmediata};
        \node[usuario, text width=4.4cm, anchor=north east] at (7.4,7.2)
        {¿Y cuál es la nota de corte?};
        \node[asistente, text width=5.8cm, anchor=north west] (r4) at (0.3,6.2)
        {Lo siento, no tengo información sobre la nota de corte del Grado en
          Ingeniería Informática. Mi contexto solo incluye información sobre las
          asignaturas y su contenido.};
        \node[pie, anchor=north west, xshift=0.1cm, yshift=-0.1cm] at (r4.south west) {09:04 \quad 23,0 s};
        \node[nota, text width=5.6cm, anchor=north west] at (0.4,3.5)
        {Ni una ni otra se pintan como error: la primera no llegó a consultar al
          modelo y la segunda es la garantía de dominio funcionando.};
        \node[entrada, text width=6.5cm, anchor=north west] at (0.3,1.5)
        {Escribe tu pregunta\ldots};
        \node[rotulo] at (3.8,-0.2) {3 · Inmediata y sin información};
      \end{scope}
    \end{tikzpicture}%
  }
   \caption[Estados de interacción de la pantalla de conversación]{Comportamiento de la interfaz ante respuestas elaboradas, tiempos de espera, respuestas inmediatas y rechazo de consultas. Fuente: Elaboración propia.}
  \label{fig:estados_chat}
\end{figure}
